
\medskip
\noindent\textbf{Proposition 2.4.1.}
\phantomsection
\label{sga3:VIIB:2.4.1}
{\itshape Let \(f:G\to K\) be a morphism of formal \(k\)-groups. If
\[
  H=\ker(f)
\]
is topologically flat over \(k\), the homomorphism
\[
  f':G/H\longrightarrow K
\]
induced by \(f\) is a monomorphism.}

This is a consequence of the results of Exposé IV,\footnote{%
\label{sga3:VIIB:nde:93}See \SGARefLink{sga3:iv:remark:5:2:6}{Exposé IV, 5.2.6}.} but we give a direct
proof. Let \(T\) be a formal variety of finite length over \(k\), and
let \(t\in(G/H)(T)\) be such that \(f'\circ t\) is the identity
element of \(K(T)\). We must show that \(t\) is the identity element
of \((G/H)(T)\). Let
\[
  p:G\longrightarrow G/H
\]
be the projection, and let
\[
  X=T\times_{G/H}G.
\]

By \SGARefLink{sga3:VIIB:2.4}{2.4}, \(p\) is surjective and topologically flat. The same is
therefore true of
\[
  \operatorname{pr}_1:X\longrightarrow T.
\]
By \SGARefLink{sga3:VIIB:1.3.1}{Proposition 1.3.1}, \(\operatorname{pr}_1\) is consequently an
epimorphism, so it is enough to show that
\(t\circ\operatorname{pr}_1\) is the identity element of
\((G/H)(X)\).

Let \(\operatorname{pr}_2:X\to G\) be the second projection. We have
\[
  t\circ\operatorname{pr}_1
  =
  p\circ\operatorname{pr}_2,
\]
and hence
\[
  1
  =
  f'\circ t\circ\operatorname{pr}_1
  =
  f'\circ p\circ\operatorname{pr}_2
  =
  f\circ\operatorname{pr}_2.
\]
The exact sequence
\begin{center}
  \phantomsection
  \label{sga3:VIIB:2.4.1:exact-sequence}
  \(
1\longrightarrow H\longrightarrow G\xrightarrow{\ f\ }K
\)

\end{center}
therefore shows that \(\operatorname{pr}_2\) factors through \(H\).
Thus \(p\circ\operatorname{pr}_2\) is the zero morphism. Since
\[
  p\circ\operatorname{pr}_2
  =
  t\circ\operatorname{pr}_1,
\]
this proves the proposition.

The following corollary follows. Write
\(\mathcal O(G)\), \(\mathcal O(K)\), and \(\mathcal O(G/H)\) for
the affine algebras of \(G\), \(K\), and \(G/H\), respectively. We
saw in \SGARefLink{sga3:VIIB:2.4}{2.4} that \(p\) induces an injection
\[
  \mathcal O(G/H)\lhook\joinrel\longrightarrow\mathcal O(G).
\]
Moreover, since \(f':G/H\to K\) is a monomorphism, \SGARefLink{sga3:VIIB:1.3.prop}{Proposition 1.3}
shows that
\[
  \mathcal O(K)\longrightarrow\mathcal O(G/H)
\]
is surjective. Hence:

\medskip
\noindent\textbf{Corollary.}
\phantomsection
\label{sga3:VIIB:2.4.1:corollary}
{\itshape Let \(f:G\to K\) be a morphism of formal \(k\)-groups and
let \(H=\ker(f)\). If \(H\) is topologically flat over \(k\), then
\[
  \mathcal O(G/H)
\]
is the image of \(\mathcal O(K)\) in \(\mathcal O(G)\).}
